\documentclass{article}
\usepackage[utf8]{inputenc}
\usepackage{amsmath, amssymb, amsfonts}
\usepackage{geometry}
\geometry{a4paper, margin=1in}
\usepackage{hyperref}

\title{LQG Notes}
\author{Whitney}
\date{\today}

\begin{document}
\maketitle
\tableofcontents

\section{Refined Algebraic Quantization}


As we are all well aware, the whole quantization programme of QFT is based on the so-called Wightman Axioms. We can loosely summarize this procedure into four steps: fixing a background metric $\eta_{\mu\nu}$, finding its corresponding symmetric group $\mathcal{P}$, formalizing the commutation relations and finally defining the vacuum state. From this, one can construct the full Fock space, which is the space of physical states. However, if one is to generalize this procedure to General Relativity, one will immediately run into trouble for failing to fix a background. For this and many other reasons, we must search for another way to quantize General Relativity.

Luckily, GR can be formulated as a constrained Hamiltonian system with first class constraints. The quantization of such a system had been explored by Dirac and many others and is known by the name of \textbf{Refined Algebraic Quantization}. This note is written in the aim of refining my understanding on this topic.

For any system, we start with a phase space $\mathcal{M}$. It has a set of constraints $C_I$ (In the case of GR, these constraints reflects the degrees of freedom in foliation and choosing spatial cordinates on hypersurfaces.) and a Hamiltonian $\mathcal{H}$.

Here let's stop to consider what a phase space with constraints will look like. $C_I$ limits $\mathcal{M}$ to a hypersurface $\bar{\mathcal{M}}$. It is also the generator of gauge transformations that leaves $\bar{\mathcal{M}}$ invariant. Therefore if we consider a point $m \in \bar{\mathcal{M}}$, a gauge $C$ generates an orbit $[m]$ on $\bar{\mathcal{M}}$. We can now assert that the physical phase space is the equivalence class of $[m]$.

For simplicity, we choose a polarization on $\mathcal{M}$ such that $\mathcal{M}$ is a cotangent bundle, and the configuration space $\mathcal{C}$ as its base manifold: $\mathcal{M} = T^\ast \mathcal{C}$. We will also assume there exist a lagrangian $\mathcal{L}$ and vanishing poisson brackets on $\mathcal{C}$. With these assumptions, we can always consider the fiber of $\mathcal{M}$, that is the momentum space. Then we compute the poisson bracket between the momentum and a function $f(q) \in C^\infty(\mathcal{C})$: $\{p,f\}(q)$. These brackets have the same structure as a deriviative of $f$: $\{p,\cdot\}(f)(q)$. It is easy to recognize this object as a vector on configuration space $v\in V^\infty(\mathcal{C})$, therefore we denote $v_p[f] := \{p,f\}(q)$, a vector field generated by $p$. This vector field no doubt carries the physical dynamics of the system and preserves the sympletric structure. We can also consider the Lie algebra $C^\infty(\mathcal{C})\times V^\infty(\mathcal{C})$ defined by $[(f,v),(f^\prime,v^\prime)] = (v[f^\prime] - v^\prime[f], [v,v^\prime])$, choosing the preferred vector field $v_p$ gives us a subalgebra $(f, v_p)$ denoted by $\mathcal{B}$.

Equipped with this algebra $\mathcal{B}$, we can construct the corresponding Hilbert space. This step is similar to the process of constructing Fock space with Poincare group in QFT, but here we won't result in a space of physical states, rather a kinematical space.
We need a irreducible $*$-representation $\pi:\mathcal{B}\to L(\mathcal{H}_{kin})$. That is to say, $\pi$ is a linear operator on $\mathcal{H}_{kin}$ such that $*$-operations are adjoints and $\forall a,b\in \mathcal{B}$,

$$
\begin{cases}
    \pi(a)^\dagger = \pi(a^\ast),\\
    [\pi(a),\pi(b)] = \mathrm{i}\hbar\pi([a,b]).
\end{cases}
$$

As stated by the \textbf{Groenewald-van Hove Theorem}, it is impossible to faithfully represent the full poisson algebra.

\begin{quote}
One can usually describe a Hilbert space in the form $L_2(\bar{\mathcal{C}},\mathrm{d}\mu)$, where $\bar{\mathcal{C}}$ is an extension of the configuration space, and $\mathrm{d}\mu$ a measure.
\end{quote}

Even though Groenewald-van Hove Theorem limits the numbers of representations we can take, there is still plenty of reps for us to choose from. We want our reps to also represent the constraints and Hamiltonians into operators. But these corresponding operators are often unbounded and therefore not well-defined on all of $\mathcal{H}_{kin}$. We need a dense subspace which is close to Hamiltonians and constraints and can avoid infinities. We accomplish this by choosing a space of rapid decrease smooth functions $\mathcal{D}_{kin}$. By this, we mean $\mathcal{D}_{kin}$ has a finer topology than $\mathcal{H}_{kin}$, thus a topological inclusion

$$
\mathcal{D}_{kin}\hookrightarrow \mathcal{H}_{kin}.
$$

We now impose our constraints. The quantum version of the constraint equation $\hat{C}_Il = 0$ is basicly requiring $l$ to be invariant under gauge transformations. Because of this, when we compute the total probability of a physical state $l$, we should integrate $l$ along its gauge orbit:

$$
\left\|l\right\|^2 = \int_{-\infty}^{+\infty}|l|^2\mathrm{d}\mu
$$

which obviously diverges. This means physical states, namely states that solve the constraint equation cannot lie in $\mathcal{H}_{kin}$. Observe that $\hat{C}_I$ has continuous spectra; for self-adjoint operators like this, \textbf{Gel'fand-Maurin Theorem} guides us to find its eigenvectors in the algebraic dual space of a nuclear space of the Hilbert space, which is exactly $\mathcal{D}^\ast_{kin}$. This by definition is a space of linear functionals, and has pointwise divergence. Again $\mathcal{H}_{kin}$ has a finer topology then $\mathcal{D}^\ast_{kin}$. We now have the \textbf{Gel'fand triple}:

$$
\mathcal{D}_{kin}\hookrightarrow \mathcal{H}_{kin}\hookrightarrow \mathcal{D}^\ast_{kin}.
$$

We now define the dual of the constraint operator as

$$
\href{f}{\hat{C}'_I l} := l(\hat{C}_I^\dagger f)
$$

where $l \in \mathcal{D}^\ast_{kin}$ and $f \in \mathcal{D}_{kin}$. It is a natural extension of the following relation in $\mathcal{H}_{kin}$:

$$
\href{f}{\hat{C}^\prime_I l} = \langle\hat{C}_I l, f\rangle_{kin} = \langle l, \hat{C}^\dagger_If\rangle_{kin} = l(\hat{C}_I^\dagger f).
$$

Now the constraint equation becomes

$$
\href{f}{\hat{C}'_I l} = l(\hat{C}_I^\dagger f) = 0\quad \forall f\in \mathcal{D}_{kin}.
$$

elements of $\mathcal{D}^\ast_{kin}$ that solve this equation form a subspace $\mathcal{D}^\ast_{phys}\subset \mathcal{D}^\ast_{kin}$.

\begin{quote}
Quantum theory comes with anomalies. Constraints come with an algebra and a structural function $\{C_I, C_J\} = f_{IJ}^KC_K$. After quantizing, the algebra becomes $[\hat{C_I}, \hat{C_J}] = \mathrm{i}\hbar\hat{f_{IJ}^K}\hat{C_K} = \mathrm{i}\hbar\{[\hat{C}_K, \hat{f_{IJ}^K}] + \hat{f_{IJ}^K}\hat{C_K}\}$. Any physical state would now also solve the equation $[\hat{C}_K, \hat{f_{IJ}^K}]l = 0$, which itself may not be a constraint. Therefore the quantized theory may have less physical degrees of freedom than the classical theory.
\end{quote}

As we've discussed, $\mathcal{H}_{kin} \cap\mathcal{D}^\ast_{phys} = \emptyset$, which is annoying because we can't equip $\mathcal{D}^\ast_{phys}$ with a scalar product. To resolve this, we introduce \textbf{Dirac Observables}. A \emph{strong} dirac observable is an operator on $\mathcal{H}_{kin}$ that commutes with all constraints and are densely defined on $\mathcal{H}_{kin}$. A \emph{weak} dirac observable only commutes with the constraint on the constraint hypersurface: $[\hat{O},\hat{C}_{I}]_{|C_{J}=0}=0$. These observables also have their natural dual $\hat{O}^prime$ on $\mathcal{D}^\ast_{kin}$, then a weak dirac observable is an operator that leaves the physical space invariant: $\hat{O}^\prime\mathcal{D}^\ast_{phys}\subset \mathcal{D}^\ast_{phys}$.

Now we can finally define a inner product on a subset $\mathcal{H}_{phys}\subset\mathcal{D}^\ast_{phys}$ as a positive definite sesquilinear form with respect to which $\hat{O}^\prime$ becomes self-adjoint operators, that is $\hat{O}^\prime = (\hat{O}^\prime)^\ast$. We also have $[\hat{O}_1^\prime, \hat{O}_2^\prime] = ([\hat{O}_1, \hat{O}_2])^\prime$. The commutator on the RHS is defined on $\mathcal{H}_{kin}$, but the LHS is defined on $\mathcal{H}_{phys}$. This is a nice relation because the commutation relations on $\mathcal{H}_{kin}$ are automatically tranferred to $\mathcal{H}_{phys}$. Here we still need a dense subspace $\mathcal{D}_{phys}\subset\mathcal{H}_{phys}$ where we have well defined operators for reasons explained above. We have finally resulted in a second Gel'fand triple:

$$
\mathcal{D}_{phys}\hookrightarrow \mathcal{H}_{phys}\hookrightarrow \mathcal{D}^\ast_{phys}.
$$

This is where the physics happen.

\begin{quote}
The logic here is that when we take our dirac observables and force them to be self-adjoint, we hope these restrictions are powerful enough to limit the form of the inner product.
\end{quote}

In very fortunate cases where constraints are mutually self-adjointing operators, namely $[\hat{C}_I, \hat{C}_J] = 0$, all constraints share the same eigenstate, therfore by spectrum theorem we can decompose $\mathcal{H}_{kin}$ into

$$
\mathcal{H}_{kin} = \int^\oplus_S\mathrm{d}\mu(\lambda)\mathcal{H}_\lambda
$$

over the spectrum $S$ of the constraint algebra. Then it's obvious that $\mathcal{H}_{phys} = \mathcal{H}_0$.

\begin{quote}
The whole process is basicly $\mathcal{H}_{kin}\to \mathcal{D}_{kin}\to \mathcal{D}^\ast_{kin} \to \mathcal{D}^\ast_{phys}\to \mathcal{H}_{phys}\to \mathcal{D}_{phys}$.
\end{quote}

The last step is to make sure this quantum theory actually has the correct classical limit. In other words, we have to make sure that

$$
\
\left|\frac{\langle\psi_{[m]},\hat{O}^\prime\psi_{[m]}\rangle_{phys}}{O(m)} - 1\right| \ll 1, \qquad \left|\frac{\langle\psi_{[m]},(\hat{O}^\prime)^2\psi_{[m]}\rangle_{phys}}{(\langle\psi_{[m]},\hat{O}^\prime\psi_{[m]}\rangle_{phys})^2} - 1\right| \ll 1.
$$

This means the dirac observables have correct expectation values and their relative fluctuations are small. Only then we are sure we have constructed a quantized theory that admits the correct classical limit.

In the next post, we will discuss how these abstract mathematical procedures can be used to really QUANTIZE a theory with constraints.

\textbf{Reference}:

\begin{itemize}
\item \href{https://arxiv.org/abs/gr-qc/0210094}{arXiv:gr-qc/0210094v1}
\end{itemize}



\section{Some Trivial Examples of RAQ}


In this post we put forward four examples of quantizing systems with constraints using the RAQ programme. In each case our final goal is to derive the dirac observables and the physical hilbert space.

\hrulefill

\subsection{Momentum Constraint}

Consider the phase space $\mathcal{M} = T^\ast \mathbb{R}^2$ with poisson brackets among $q^a$ and $p_a$. Our constraint is $C = p_1$, and we choose $\mathcal{H}_{kin} = L_2(\mathbb{R}^2, \mathrm{d}^2x)$, $\mathcal{D}_{kin} = \mathcal{S}(\mathbb{R}^2)$ and $\mathcal{D}^\ast_{kin} = \mathcal{S}^\ast(\mathbb{R}^2)$.

By the constraint $C = p_1$, we have the constraint equation $l(\hat{C}^\dagger\phi) = 0$, where $l \in \mathcal{D}^\ast_{kin}$ and $\phi \in \mathcal{D}_{kin}$. Working in the momentum representation, we can write the constraint equation as

$$
l(\hat{p}_1\phi(\hat{p}_1, \hat{p_2})) = 0
$$

where we can cleanly see that a function of the form $l = \delta(p_1)f(p_2)$ solves the equation. We construct the new inner product as $\langle l_{f_1}, l_{f_2}\rangle := l_1[\phi_2]$, therefore we have

$$
l_1[\phi_2] = \int \int \delta(p_1) f_1^*(p_2) \phi_2(p_1, p_2) \mathrm{d}p_1 \mathrm{d}p_2 = \int f_1^*(p_2) \phi_2(p_2) \mathrm{d}p_2.
$$

which is a perfectly well-defined inner product, therefore $\mathcal{H}_{phys} = L_2(\mathbb{R},\mathrm{d}x_2)$.

From the constraint it's also trivial to find the dirac observables to be $q^2, p_2$ because these operators must commute with the constraint, which in this case is $p_1$.

\hrulefill

\subsection{Angular Momentum Constraint}

We consider the phase space to be $\mathcal{M} = T^\ast \mathbb{R}^3$ with poisson brackets among $q^a$ and $p_a$. Our constraints are $C_a = \epsilon_{abc}x_bp_c$, and we choose $\mathcal{H}_{kin} = L_2(\mathbb{R}^3, \mathrm{d}^3x)$, $\mathcal{D}_{kin} = \mathcal{S}(\mathbb{R}^3)$ and $\mathcal{D}^\ast_{kin} = \mathcal{S}^\ast(\mathbb{R}^3)$.

The constraint algebra is simply the angular momentum algebra, so we can identify $\hat{C}_a \equiv \hat{L}_a$. The constraint equation is now $\hat{L}_a \Psi = 0$. By expanding it into spherical harmonics, we have

$$
\Psi = \sum_{l = 0}^{\infty}\sum_{m = -l}^lR_{lm}(r)Y^m_l(\theta,\phi).
$$

The constraint equation implies $\hat{L}^2\Psi = 0$, but from QM we know that $\hat{L}^2\Psi = l(l+1)\hbar^2Y^m_l$, thus $l = 0$. Our physical states are then $\Psi = \phi(r)$, which can be perfectly normalized in $\mathcal{H}_{kin}$ because $SO(3)$ is compact and the integral on the group manifold gives a finite result. Therefore, we have

$$
\mathcal{H}_{phys} = L_2(\mathbb{R},r^2\mathrm{d}r) \subset \mathcal{H}_{kin}.
$$

Now we find the dirac observables. We'll have to find operators that are self-adjoint under the physical inner product induced by $\mathcal{H}_{kin}$. Notice that if we define $\hat{p} = -\mathrm{i}\hbar r^{-1}\frac{\mathrm{d}}{\mathrm{d}r}r$, such a inner product can be written as

$$
\begin{aligned}
    \langle\phi, \hat{p}\psi\rangle = \int \phi^\ast\hat{p}\psi r^2\mathrm{d}r
    &=  \int \phi^\ast\left(-\mathrm{i}\hbar \frac{1}{r}\frac{\mathrm{d}}{\mathrm{d}r}\left(r\psi\right) \right)r^2\mathrm{d}r\\
    &=  -\mathrm{i}\hbar \int r\phi^\ast\left(\partial_r\left(r\psi\right) \right)\mathrm{d}r\\
    &=  \left(-\mathrm{i}\hbar(r\phi^\ast)(r\phi)\right)\Big|^{+\infty}_{-\infty} + \mathrm{i}\hbar \int \partial_r(r\phi^\ast)\left(r\psi\right)\mathrm{d}r\\
    &=  \int\left(\hat{p}\phi\right)^\ast\psi r^2\mathrm{d}r = \langle\hat{p}\phi, \psi\rangle.
\end{aligned}
$$

Similarly the operator $\hat{r} := r$ is also a dirac observable, completing the programme.

\hrulefill

\subsection{Relativistic Particle}

From the form of the Lagrangian

$$
L = -m\sqrt{-\eta_{\mu\nu} \dot{q}^\mu\dot{q}^\nu}
$$

and the definition of conjugate momenta $p_\mu = \partial L/\partial \dot{q}^\mu$, we can surprisingly see

$$
\frac{\partial L}{\partial \dot{q}^\mu} = m\frac{\dot{q}_\mu}{\sqrt{-\eta_{\mu\nu} \dot{q}^\mu\dot{q}^\nu}} = p_\mu \implies \eta^{\mu\nu}p_\mu p_\nu + m^2 = 0.
$$

which is the mass-shell constraint. In this case, the constraint came directly from the lagrangian. We shall later see this is the result of the disappearing Hamiltonian.

Next we prove the action $S = \int \mathrm{d}t L$ to be $\mathrm{Diff}(\mathbb{R})$ invariant. Consider the reparameterization $t\mapsto \varphi(t)$. With this new time variable, the action is now

$$
S^\prime = \int -m\sqrt{-\eta_{\mu\nu}\frac{\mathrm{d}\dot{q}^\mu}{\mathrm{d}\varphi}\frac{\mathrm{d}\dot{q}^\nu}{\mathrm{d}\varphi}}\mathrm{d}\varphi = \int -m\sqrt{-\eta_{\mu\nu}\frac{\mathrm{d}\dot{q}^\mu}{\mathrm{d}t}\frac{\mathrm{d}\dot{q}^\nu}{\mathrm{d}t}}\dot{\varphi}\dot{\varphi}^{-1}\mathrm{d}t = S.
$$

Under reparameterization, although $\dot{q}^\mu$ is variant, $p_\mu$ stayed constant. This means the mapping from velocity space to momentum space is not one-to-one, namely every momentum correpsonds to a ray in velocity space. Therefore we need to have some kind of constraint on $\mathcal{M}$, limiting the physical momentum to a hypersurface, here the mass-shell constraint.

Next we calculate the hamiltonian, which is trivially vanishing:

$$
H = p_\mu \dot{q}^\mu - L = 0.
$$

This means in this system, evolution is simply a gauge transformation and there are no dynamics.

We consider the phase space to be $\mathcal{M} = T^\ast \mathbb{R}^{D+1}$ with poisson brackets among $q^a$ and $p_a$. Our constraint is the mass-shell constraint, and we choose $\mathcal{H}_{kin} = L_2(\mathbb{R}^{D+1}, \mathrm{d}^{D+1}x)$, $\mathcal{D}_{kin} = \mathcal{S}(\mathbb{R}^{D+1})$ and $\mathcal{D}^\ast_{kin} = \mathcal{S}^\ast(\mathbb{R}^{D+1})$. This system is rather similar to the momentum one, so we work in the momentum representation again to get

$$
l(\hat{C}\Psi) = 0 \implies l((-\partial^2 + m^2)\Psi) = 0 \implies l((-p_0^2 + p^2 + m^2)\Psi) = 0.
$$

We assume the solution is of the form $l_f = \delta(-p_0^2 + p^2 + m^2)f(p_0,p_i)$. Applying the same physical inner product as in 1, we can get

$$
\langle l_{f_1}, l_{f_2}\rangle = \int l^\ast_{f_1}f_2(p_i)\mathrm{d}^Dp.
$$

that is to say the physical hilbert space is $\mathcal{H}_{phys} := L_2(\mathbb{R}^2,\mathrm{d}^Dp)$. The dirac observables are $Q^a = q^a - q^0\frac{p_a}{\sqrt{m^2+p^2}}$. It is easy to verify that for every $q^0$, $[Q^a,C] = 0$. This is another sign that coordinate time in relativistic systems are gauge, and for simplicity reasons we set $q^0 = 0$.

\hrulefill

\subsection{Maxwell Theory}

As we all know, the lagrangian for a free maxwell theory is $\mathcal{L} = -1/4F_{\mu\nu}F^{\mu\nu}$. We define $E^a = \dot{A}_a - \partial_a A_0$, and $B^a = \epsilon^{abc}\partial_b A_c$. The lagrangian is therefore

$$
L = \int \mathrm{d}^3x\frac{1}{2}\left(E^aE_a - B^aB_a\right)
$$

the generalized coordinates in this case are $A^\mu$. If we compute the conjugate momenta of $A^\mu$, we find

$$
\begin{cases}
    \frac{\partial L}{\partial \dot{A}_0} = 0 =: \pi^0,\\
    \frac{\partial L}{\partial \dot{A}_a} = E^a =: \pi^a.
\end{cases}
$$

and the hamiltonian to be

$$
H = \int\mathrm{d}^3x\left(\pi^0 \dot{A}_0 + \pi^a\dot{A}_a - \mathcal{L}\right) = \int\mathrm{d}^3x\frac{1}{2}(\pi^2 + B^2) + \pi^a\partial_a A_0.
$$

For the term $\pi^a\partial_a A_0$, we perform integration by parts: $\pi^a\partial_a A_0 = -A_0\partial_a\pi^a$. The time evolution of $\pi^0$ should be 0, which leads to the poisson bracket $\{\pi^0, H\} = -\delta H/\delta A_0 = \partial_a\pi^a = \partial_a E^a = 0$. But this is exactly the Gauss constraint! $A_0$ has lost all it's dynamics under this constraint, but only a lagrange multiplier.

Now we demonstrate that the Gauss constraint does indeed generate the $U(1)$ gauge transformation. We consider the integration of $\partial_a E^a$ mutiplied by some function $\lambda(x)$:

$$
C(\lambda) = \int \mathrm{d}^3x \lambda(x) \partial_a E^a = - \int d^3x (\partial_a \lambda) E^a
$$

Varying $A$ gives

$$
\delta A(x) = \{A_a(x), C(\lambda)\} = -\partial_a\lambda(x),
$$

which is the standard $U(1)$ gauge transformation. Similarly, we can see that $E$ is gauge invariant.

In order to do quantization, we first apply the Helmholtz decomposition on $A_a$ and $E^a$:

$$
\begin{cases}
    A_a = A^L_a + A^T_a,\\
    E^a = (E^L)^a + (E^T)^a.
\end{cases}
$$

where $\nabla\cdot A^T = 0, A^L_a = \partial_a\phi$ and $\nabla\cdot E^T = 0, (E^L)^a = \partial_a\psi$. Now the Gauss constraint becomes $C = \partial_a E^a = \partial_a (E^L)^a = \nabla^2\psi = 0$. Therefore, the gauge condition suggests that the dirac observables are the transversal parts of $A$ and $E$.

Now we do the quantization. Consider the three corresponding massless scalar fields of $A_a$. $\mathcal{H}_{kin}$ is the Fock space generated by $\hat{z}^\dagger_a(k)$ and $\hat{z}_a(k)$, the creation and annihilation operators respectively. Under fourier transform, the constraint can be written as $\hat{C} = k_aE^a(k)$, which leads to $k^a\hat{z}_a(k)|\Psi\rangle_{phys} = 0$. Next we introduce polarized vectors $e_a(k)$. But because by the gauss constraint $E^a$ only consist of components parallel to $k$, the constraint equation is simply $\hat{z}_3(k)|\Psi\rangle_{phys} = 0$. This means physical states are states without longitudonal excitations, therefore $\mathcal{H}_{phys}$ is the Fock space spanned by $\hat{z}^\dagger_1(k)$ and $\hat{z}^\dagger_2(k)$. This is why there are only two tranverse polarization states of a photon.

\hrulefill

In the next part, we will use the ADM formulation to construct the phase space of General Relativity.


\section{The Connection Formulation of General Relativity}


It has long been known that general relativity can be formulated by applying a constraint on the BF theory. Plebanski successfully obtained the Riemannian action:

$$
S[B, \omega, \lambda] = \frac{1}{\kappa}\int_\mathcal{M}\left[\left(^\star\!B^{IJ} + \frac{1}{\gamma}B^{IJ}\right)\wedge F_{IJ}(\omega) + \lambda_{IJKL} B^{IJ}\wedge B^{KL}\right].
$$

This can be thought of as a $\mathrm{Spin}(4)$ BF theory and a constraint imposed by the lagrangian multiplier $\lambda$. It depends on a $\mathfrak{so}(4)$ connection $\omega$ and a Lie-algebra valued 2-form $B$.

The symmetry of $B$ and $\wedge$ imposes a constraint on $\lambda$ as well. Since switching between the indices of $B^{IJ}$ is antisymmetric and that the operation $\wedge$ is symmetric, we have the following relation:

$$
\lambda_{IJKL} = -\lambda_{JIKL} = -\lambda_{IJLK} = \lambda_{KLIJ}.
$$

This restrains the dimension of $\lambda$ down to 21. However, if $\lambda_{IJKL}\propto \epsilon_{IJKL}$ the lagrangian multiplier will give $\epsilon_{IJKL}B^{IJ}\wedge B^{KL}$, which implies a degenerate metric. Therefore $\epsilon^{IJKL}\lambda_{IJKL} = 0$, and that $\mathrm{dim}\lambda = 20$.

To actually calculate the constraint given by the lagrangian multiplier, we introduce

$$
\lambda_{IJKL} = \Lambda_{IJKL} - \frac{1}{4!}\Lambda_{MNOP}\epsilon^{MNOP}\epsilon_{IJKL}.
$$

We make this substitution because if we assume $\Lambda$ to be the unconstrained lagrangian multiplier and

$$
\Lambda_{[IJKL]} = c\cdot\epsilon_{IJKL}
$$

we can multiply both sides with $\epsilon$ and have

$$
\Lambda_{[IJKL]}\epsilon^{IJKL} = c\cdot\epsilon^2\implies c = \frac{1}{4!}\Lambda_{[MNOP]}\epsilon^{MNOP}.
$$

This gives the derivative of the action with respect to $\Lambda$:

$$
\frac{\delta S}{\delta \Lambda_{IJKL}} = B^{IJ}\wedge B^{KL} - \frac{1}{4!}\epsilon^{IJKL}\epsilon_{MNOP}B^{MN}\wedge B^{OP} = 0.
$$

which implies

$$
\epsilon^{\mu\nu\rho\sigma} B^{IJ}_{\mu\nu} B^{KL}_{\rho\sigma} = \frac{1}{4!} \epsilon^{IJKL} \epsilon_{MNOP} B^{MN}_{\mu\nu} B^{OP}_{\rho\sigma} \epsilon^{\mu\nu\rho\sigma}.
$$

This equation constraints the degrees of freedom of $B$, limiting its dimension down to $16 = 4\times 4$. So if we define:

$$
e := \frac{1}{4!} \epsilon_{MNOP} B^{MN}_{\mu\nu} B^{OP}_{\rho\sigma} \epsilon^{\mu\nu\rho\sigma}
$$

It behaves exactly as the local tetrad. The solution of $B$ would be

$$
B = \pm ^\star(e\wedge e) \quad\text{and}\quad \pm e\wedge e.
$$

Here $e$ is the new variable, and the term $\lambda B\wedge B$ in the action is simply the constraint. So if we put $B = \pm^\star(e\wedge e)$ back into the action $S$, we obtain

$$
S[e, \omega] = \frac{1}{\kappa}\int_\mathcal{M} \operatorname{Tr}\left[\left(^\star(e\wedge e) + \frac{1}{\gamma}e\wedge e\right) \wedge F(\omega)\right].
$$

\subsection{Canonical Analysis}

We perform the 3+1 decomposition. We use labels $a,b = 1,2,3$ as spacial indices on $\Sigma_t$. Introduce

$$
\begin{aligned}
\Pi_{ab}^{IJ} &= ^\star B_{ab}^{IJ} + \frac{1}{\gamma}B_{ab}^{IJ}\implies \text{pure spacial; lives on }\Sigma_t\\
H_a^{IJ} &= ^\star B_{0a}^{IJ} + \frac{1}{\gamma}B_{0a}^{IJ}\implies \text{space-time}.
\end{aligned}
$$

Also, we decompose the curvature form $F_{\mu\nu}$:

$$
\begin{aligned}
F_{ab} &= \partial_{[a}\omega_{b]} + [\omega_a, \omega_b]\\
F_{0a} &= \partial_0\omega_a - \partial_a\omega_0 + [\omega_0, \omega_a] = \dot{\omega}_a - (\partial_a\omega_0 - [\omega_a, \omega_0]) =: \dot{\omega}_a - \mathcal{D}_a\omega_0.
\end{aligned}
$$

The action is originally:

$$
S[\Pi, H, \omega, \lambda] = \frac{1}{\kappa}\int_\mathcal{M}\left[\left(^\star B^{IJ} + \frac{1}{\gamma}B^{IJ}\right)\wedge F_{IJ} + \lambda_{IJKL} B^{IJ}\wedge B^{KL}\right]
$$

Substituting the 3+1 decomposed variables:

$$
S = \frac{1}{\kappa}\int_\mathcal{M} \left[(\Pi_{ab}^{IJ}\wedge F_{IJ,0a}) + (H_{a}^{IJ}\wedge F_{IJ,ab}) + \lambda_{IJKL}B_{0a}^{IJ}\wedge B_{ab}^{KL}\right]
$$

After some algebra, this leads to:

$$
\begin{aligned}
S
&= \frac{1}{\kappa}\int_\mathcal{M}\left(\Pi_{ab}^{IJ}\wedge \dot{\omega}_{IJ,c} - \Pi_{ab}^{IJ}\wedge \mathcal{D}_a \omega_{0,IJ} + H_a^{IJ}\wedge F_{IJ,ab} \right. \\
&+  \left.\lambda_{IJKL}\left\{\frac{\gamma^2}{\gamma^2-1}(\gamma^\star H_{a}^{IJ} - H_{a}^{IJ}) \wedge \frac{\gamma}{\gamma^2-1}(\gamma\Pi_{bc}^{KL} - ^\star\Pi_{bc}^{KL})\right\}\right).
\end{aligned}
$$

By integrating by parts, the second term gives $\omega_0 \mathcal{D}_\omega\wedge \Pi$, and if we redefine $\lambda' = \frac{\gamma^2}{(\gamma^2-1)^2}\lambda$, the action is:

$$
S = \int dt \int_{\Sigma_t} \operatorname{Tr} \left[ \Pi \wedge \dot{\omega} + \omega_0 \mathcal{D}_\omega \wedge \Pi + H \wedge F + \lambda' (\gamma ^\star H - H) \wedge (\gamma ^\star \Pi - \Pi) \right].
$$

The canonical pair is $(\Pi^{IJ}_{ab}, \omega^{KL}_c)$. If we introduce

$$
\begin{aligned}
\mathbb{P}(\gamma, \pm)^i_{ab} &:= \frac{1}{4} \epsilon^i_{\ jk} \Pi^{jk}_{ab} \pm \frac{1}{2\gamma} \Pi^{0i}_{ab} \\
{}^{\pm}\mathbb{A}_c^i &:= \frac{1}{2} \epsilon^i_{\ mn} \omega^{mn}_c \pm \gamma \omega^{0i}_c \\
\mathbb{H}(\gamma, \pm)^i_a &:= \frac{1}{4} \epsilon^i_{\ jk} H^{jk}_a \pm \frac{1}{2\gamma} H^{0i}_a
\end{aligned}
$$

where $\mathbb{A}$ is a function of $\omega$, its time derivative corresponds to

$$
\partial_0 {}^{\pm}\mathbb{A}_c^i = \frac{1}{2} \epsilon^i_{\ mn} \partial_0 \omega^{mn}_c \pm \gamma \partial_0 \omega^{0i}_c = \frac{1}{2} \epsilon^i_{\ mn} \dot{\omega}^{mn}_c \pm \gamma \dot{\omega}^{0i}_c =: {}^{\pm}\dot{\mathbb{A}}_c^i.
$$

Also notice that

$$
\mathbb{P}(\gamma, +)^i_{ab} = \frac{1}{4} \epsilon_{ijk} \Pi^{jk}_{ab} + \frac{1}{2\gamma} \Pi^{0}_{ab, i},
$$

$$
\mathbb{P}(\gamma, -)^i_{ab} = \frac{1}{4} \epsilon_{ijk} \Pi^{jk}_{ab} - \frac{1}{2\gamma} \Pi^{0}_{ab, i}.
$$

These simplify the canonical trace term $\operatorname{Tr}[\Pi\wedge\dot{\omega}]$:

$$
\begin{aligned}
\mathbb{P}^{+} \wedge {}^{+}\dot{\mathbb{A}} + \mathbb{P}^{-} \wedge {}^{-}\dot{\mathbb{A}}
&= (\Pi + \Pi^0) \wedge (\dot{\omega} + \dot{\omega}^0) + (\Pi - \Pi^0) \wedge (\dot{\omega} - \dot{\omega}^0) \\
&= 2 \Pi \wedge \dot{\omega} + 2 \Pi^0 \wedge \dot{\omega}^0 \\
&= \frac{1}{2} \epsilon_{ijk} \Pi^{jk} \wedge \frac{1}{2} \epsilon^i_{\, mn} \dot{\omega}^{mn} + \frac{1}{\gamma} \Pi^{0}_i \wedge \gamma \dot{\omega}^{i0} \\
&= \frac{1}{2} (\delta_{jm}\delta_{kn} - \delta_{jn}\delta_{km}) \Pi^{jk} \wedge \dot{\omega}^{mn} + \Pi^{0}_i \wedge \dot{\omega}^{i0} \\
&= \frac{1}{2} \Pi_{ij} \wedge \dot{\omega}^{ij} + \Pi^0_{i} \wedge \dot{\omega}^{i0} = \operatorname{Tr} [ \Pi \wedge \dot{\omega} ].
\end{aligned}
$$

The term $\omega_0 \mathcal{D}_\omega \wedge \Pi$ acts as a Lagrangian constraint, with $\omega_0$ serving as the multiplier. We can express the connection variables as:

$$
\frac{{}^{+}\mathbb{A} + {}^{-}\mathbb{A}}{2} = \frac{1}{2} \epsilon^i_{\ mn} \omega^{mn}_c \quad \text{and} \quad ({}^{+}\mathbb{A} - {}^{-}\mathbb{A})_j = 2\gamma \omega^{0j}_c.
$$

Consequently, we can decompose the constraint term:

$$
\begin{aligned}
\omega_0 \mathcal{D}_\omega \wedge \Pi &= {}^{+}\mathbb{A}_0 (\mathcal{D}_{{}^{+}\mathbb{A}} \mathbb{P}^{+})_i + {}^{-}\mathbb{A}_0 (\mathcal{D}_{{}^{-}\mathbb{A}} \mathbb{P}^{-})_i.
\end{aligned}
$$

Expanding the first term, we find:

$$
\begin{aligned}
{}^{+}\mathbb{A}_0 \left( d \mathbb{P}^{+} + {}^{+}\mathbb{A} \times \mathbb{P}^{+} \right) &= {}^{+}\mathbb{A}_0 \left[ d \mathbb{P}^{+} + \frac{1}{2}\epsilon_{ijk} \left( \epsilon^j_{\ mn} \omega^{mn} + \gamma \omega^{0j} \right) \mathbb{P}^{+k}\right] \\
&= {}^{+}\mathbb{A}_0 \left[ d \mathbb{P}^{+} + \frac{1}{2} \epsilon_{ijk} \epsilon^j_{\ mn} \omega^{mn} \mathbb{P}^{+k} + \gamma \epsilon_{ijk} \omega^{0j} \mathbb{P}^{+k} \right] \\
&= {}^{+}\mathbb{A}_0 \, \mathcal{D}_{\left(\frac{1}{2}\epsilon \omega\right)} \mathbb{P}^{+} + {}^{+}\mathbb{A}_0 \, \epsilon_{ijk} \frac{{}^{+}\mathbb{A} - {}^{-}\mathbb{A}}{2} \mathbb{P}^{+k} \\
&\sim \eta_i \mathbb{B}^i + N_i \mathbb{G}^i.
\end{aligned}
$$

Here we introduce new constraints $\mathbb{B}^i$ and $\mathbb{G}^i$. Hence, we can rewrite the action as:

$$
\begin{aligned}
S[\mathbb{P}(\gamma, \pm), {}^{\pm}\mathbb{A}, \lambda] &= \int dt \int_{\Sigma_t} \Big\{ \mathbb{P}^{+} \wedge {}^{+}\dot{\mathbb{A}} + \mathbb{P}^{-} \wedge {}^{-}\dot{\mathbb{A}} + N_i \mathbb{G}^i + \eta_i \mathbb{B}^i \\
&+ {}^{-}\mathbb{A}_0 \, \mathcal{D}_{{}^{-}\mathbb{A}} \mathbb{P}^{-} + H \wedge F + \lambda' (\gamma ^\star H - H) \wedge (\gamma ^\star \Pi - \Pi) \Big\}.
\end{aligned}
$$

The Lagrangian multipliers $N, \eta, \lambda'$ impose the following constraints:

$$
\begin{cases}
\mathbb{B}^i := \mathcal{D}_{\frac{{}^{+}\mathbb{A} + {}^{-}\mathbb{A}}{2}} \wedge \mathbb{P}(\gamma, +)^i \approx 0 \\
\mathbb{G}^i := \frac{1}{2\gamma} \epsilon_{ijk} ({}^{+}\mathbb{A} - {}^{-}\mathbb{A})_j \wedge \mathbb{P}(\gamma, +)_k \approx 0 \\
(\gamma^\star H - H)\wedge(\gamma^\star \Pi - \Pi) = 0 \implies \epsilon_{\mu\nu\rho\sigma} B^{IJ}_{\mu\nu} B^{KL}_{\rho\sigma} - e\epsilon^{IJKL} = 0.
\end{cases}
$$

The index pair $[IJ, KL]$ must be decomposed. We separate it into three cases:

$$
\begin{cases}
IJKL = 0i0j\\
IJKL = ijkl\\
IJKL = 0ijk.
\end{cases}
$$

For case 1 ($IJKL = 0i0j$):

$$
\begin{gathered}
B^{0i} \wedge B^{0j} = 0 \\
\implies C^{0i}_H \wedge C^{0j}_\Pi + (i \leftrightarrow j) = 0 \\
\implies (\gamma ^\star H^{0i} - H^{0i}) \wedge (\gamma ^\star \Pi^{0j} - \Pi^{0j}) + (i \leftrightarrow j) = 0 \\
\implies \left(\frac{\gamma}{2}\epsilon^i_{\ jk}H^{jk} - H^{0i}\right) \wedge \left(\frac{\gamma}{2}\epsilon^j_{\ kl}\Pi^{kl} - \Pi^{0j}\right) + (i \leftrightarrow j) = 0.
\end{gathered}
$$

For case 2 ($IJKL = ijkl$):

$$
\begin{gathered}
B^{ij} \wedge B^{kl} = 0 \\
\implies C^{ij}_H \wedge C^{kl}_\Pi + (ij \leftrightarrow kl) = 0 \\
\implies (\gamma ^\star H^{ij} - H^{ij}) \wedge (\gamma ^\star \Pi^{kl} - \Pi^{kl}) + (ij \leftrightarrow kl) = 0 \\
\implies \epsilon^{(p}_{\ \ \  ij}\left(\frac{\gamma}{2}\epsilon^{ij}_{\ \ m} H^{0m} - H^{ij}\right) \wedge \epsilon^{q)}_{\ \ \ kl}\left(\frac{\gamma}{2}\epsilon^{kl}_{\ \ n} \Pi^{0n} - \Pi^{kl}\right) = 0.
\end{gathered}
$$

For case 3 ($IJKL = 0ijk$):

$$
\begin{gathered}
B^{0i} \wedge B^{jk} = \nu \epsilon^{0ijk} \\
\implies C^{0i}_H \wedge C^{jk}_\Pi + (0i \leftrightarrow jk) = \nu \epsilon^{0ijk} \\
\implies (\gamma ^\star H^{0i} - H^{0i}) \wedge (\gamma ^\star \Pi^{jk} - \Pi^{jk}) + (0i \leftrightarrow jk) = \nu \epsilon^{0ijk} \\
\begin{aligned}
\implies &\left(\frac{\gamma}{2}\epsilon^i_{\ mn} H^{mn} - H^{0i}\right) \wedge \left(\frac{\gamma}{2}\epsilon^{jk}_{\ \ l} \Pi^{0l} - \Pi^{jk}\right) +\\
&\left(\frac{\gamma}{2}\epsilon^i_{\ mn}\Pi^{mn} - \Pi^{0i}\right) \wedge \left(\frac{\gamma}{2}\epsilon^{jk}_{\ \ l}H^{0l} - H^{jk}\right) - \nu \epsilon^{0ijk}= 0.
\end{aligned}
\end{gathered}
$$

Recall that we've defined

$$
\begin{aligned}
\mathbb{H}(\gamma, \pm)^i_a &:= \frac{1}{4} \epsilon^i_{\ jk} H^{jk}_a \pm \frac{1}{2\gamma} H^{0i}_a \\
\mathbb{P}(\gamma, \pm)^i_{ab} &:= \frac{1}{4} \epsilon^i_{\ jk} \Pi^{jk}_{ab} \pm \frac{1}{2\gamma} \Pi^{0i}_{ab}.
\end{aligned}
$$

Therefore we can rewrite the constraints as:

$$
\begin{aligned}
&\mathbb{B}^i = \mathcal{D}_{({}^{+}\mathbb{A} + {}^{-}\mathbb{A})/2} \wedge \mathbb{P}(\gamma, +)^i \approx 0 \\
&\mathbb{G}^i = \frac{1}{2\gamma} \epsilon^{ijk} ({}^{+}\mathbb{A} - {}^{-}\mathbb{A})_j \wedge \mathbb{P}(\gamma, +)_k \approx 0 \\
&\mathbb{I}^{ij} = \mathbb{H}(\gamma, -)^{(i} \wedge \mathbb{P}(\gamma, -)^{j)} \approx 0 \\
&\mathbb{II}^{ij} = \mathbb{H}\left(\frac{1}{\gamma}, -\right)^{(i} \wedge \mathbb{P}\left(\frac{1}{\gamma}, -\right)^{j)} \approx 0 \\
&\mathbb{III}^{ij} = \mathbb{H}(\gamma, -)^{i} \wedge \mathbb{P}(\gamma, +)^{j} + \mathbb{H}\left(\frac{1}{\gamma}, -\right)^{i} \wedge \mathbb{P}(\gamma, -)^{j} - \frac{1}{3} \delta^{ij} \operatorname{Tr}[\dots] = 0
\end{aligned}
$$

The last three constraints are derived from the original lagrangian multiplier $\lambda$, therefore include 20 constraints in total and are simplicity constraints. In the 3+1 setting, by choosing the time gauge to be $n^I = (1, 0, 0, 0)$, $e^I_\mu$ is

$$
\begin{cases}
e^0_0 = N & (\text{lapse}) \\
e^0_a = 0 & (\text{pure spacial, no time-like component}) \\
e^i_0 = N^i & (\text{shift}) \\
e^i_a = \text{standard 3-d tetrad}
\end{cases}
$$

(Indices: $I, J \in 0, 1, 2, 3$ local; $\mu, \nu \in 0, 1, 2, 3$ global; $a,b \in 1, 2, 3$ global spacial; $i,j \in 1, 2, 3$ local spacial).
We consider the solution $B = e \wedge e \implies B^{IJ}_{\mu\nu} = e^I_\mu e^J_\nu - e^J_\mu e^I_\nu$.

Globally spacial components:

$$
\begin{aligned}
B^{0i}_{ab} &= e^0_a e^i_b - e^0_b e^i_a = 0 \implies ^\star B^{0i}_{ab} = \frac{1}{2} \epsilon_{\ \ \ jk}^{0i} B^{jk}_{ab} = \frac{1}{2} \epsilon^{0i}_{\ \ \ jk} (e^j_a e^k_b - e^j_b e^k_a) = \epsilon^{i}_{\ \ jk} e^j_a e^k_b \\
B^{jk}_{ab} &= e^j_a e^k_b - e^k_a e^j_b = 2 e^{[j}_a e^{k]}_b \implies ^\star B^{jk}_{ab} = \frac{1}{2} \epsilon^{jk}_{\ \ \ 0i} B^{0i}_{ab} = 0.
\end{aligned}
$$

Globally timelike components:

$$
\begin{aligned}
B^{0i}_{0a} &= e^0_0 e^i_a - e^0_a e^i_0 = N e^i_a \implies ^\star B^{0i}_{0a} = \frac{1}{2} \epsilon_{\ \ \ jk}^{0i} B^{jk}_{0a} = \frac{1}{2} \epsilon^{0i}_{\ \ jk} (e^j_0 e^k_a - e^k_0 e^j_a) = \epsilon^{i}_{\ \ jk} N^j e^k_a \\
B^{jk}_{0a} &= e^j_0 e^k_a - e^k_0 e^j_a = 2 N^{[j} e^{k]}_a \implies ^\star B^{jk}_{0a} = \frac{1}{2} \epsilon^{jk}_{\ \ \ 0i} B^{0i}_{0a} = \frac{1}{2} \epsilon^{jk}_{\ \ \ 0i} N e^i_a = N \epsilon^{jk}_{\ \ \ 0i} e^{i}_{\ \ a}.
\end{aligned}
$$

Now we calculate the corresponding $\Pi$ \& $H$:

$$
\begin{aligned}
\Pi^{0i}_{ab} &= {}^\star B^{0i}_{ab} + \frac{1}{\gamma} B^{0i}_{ab} = \epsilon^{i}_{\ \ jk} e^{j}_{a} e^{k}_{b} \\
\Pi^{jk}_{ab} &= {}^\star B^{jk}_{ab} + \frac{1}{\gamma} B^{jk}_{ab} = \frac{2}{\gamma} e^{[j}_{a} e^{k]}_{b} \\
H^{0i}_{a} &= {}^\star B^{0i}_{0a} + \frac{1}{\gamma} B^{0i}_{0a} = \epsilon^{i}_{\ \ jk} N^{j} e^{k}_{a} + \frac{N}{\gamma} e^{i}_{a} \\
H^{jk}_{a} &= {}^\star B^{jk}_{0a} + \frac{1}{\gamma} B^{jk}_{0a} = N \epsilon_{\ \ i}^{jk} e^{i}_{a} + \frac{2}{\gamma} N^{[j} e^{k]}_{a}.
\end{aligned}
$$

Now we can show that

$$
\begin{aligned}
\mathbb{P}(\gamma, \pm)^i_{ab} &= \frac{1}{4} \epsilon^i_{\ jk} \Pi^{jk}_{ab} \pm \frac{1}{2\gamma} \Pi^{0i}_{ab} = \frac{1}{4} \epsilon^i_{\ jk} \frac{2}{\gamma} e^j_{[a} e^k_{b]} \pm \frac{1}{2\gamma} \epsilon^i_{\ jk} e^j_a e^k_b =
\begin{cases}
\mathbb{P}(\gamma, -)^i_{ab} = 0 \\
\mathbb{P}(\gamma, +)^i_{ab} = \frac{1}{\gamma} \epsilon^i_{\ jk} e^j_a e^k_b
\end{cases} \\
\mathbb{P}\left(\frac{1}{\gamma}, \pm\right)^i_{ab} &= \frac{1}{4} \epsilon^i_{\ jk} \Pi^{jk}_{ab} \pm \frac{\gamma}{2} \Pi^{0i}_{ab} = \frac{1}{2\gamma} \epsilon^i_{\ jk} e^j_a e^k_b \pm \frac{\gamma}{2} \epsilon^i_{\ jk} e^j_a e^k_b = \frac{1\pm\gamma^2}{2\gamma}\mathbb{P}(\gamma, +)^i_{ab} \\
\mathbb{H}(\gamma, -)^i_a &= \frac{1}{4} \epsilon^i_{\ jk} H^{jk}_a - \frac{1}{2\gamma} H^{0i}_a = \frac{1}{4} \epsilon^i_{\ jk} \left(N \epsilon^{jk}_{\ \ \ l} e_{a}^l + \frac{2}{\gamma} N^{[j} e^{k]}_a\right) - \frac{1}{2\gamma}\left(\epsilon^{i}_{\ \ jk} N^j e_{a}^k + \frac{N}{\gamma} e^i_a\right) \\
&= \frac{1}{2} N \delta^i_l e^l_a + \frac{1}{2\gamma} \epsilon^i_{\ jk} N^j e^k_a - \frac{1}{2\gamma} \epsilon^{i}_{\ \ jk} N^j e_{a}^k - \frac{N}{2\gamma^2}e^i_a = N \frac{\gamma^2 - 1}{2\gamma^2} e^i_a \\
\mathbb{H}\left(\frac{1}{\gamma}, -\right)^i_a &= \frac{1}{4} \epsilon^i_{\ jk} H^{jk}_a - \frac{\gamma}{2} H^{0i}_a = \frac{1}{2} N \delta^i_l e^l_a + \frac{1}{2\gamma} \epsilon^i_{\ jk} N^j e^k_a - \frac{\gamma}{2}\left(\epsilon^{i}_{\ \ jk} N^j e_{a}^k + \frac{N}{\gamma} e^i_a\right) \\
&= \frac{1-\gamma^2}{2\gamma} \epsilon^{i}_{\ \ jk} N^b e^j_b e^k_a = \frac{1 - \gamma^2}{2} N^b \mathbb{P}(\gamma, +)^i_{ba}.
\end{aligned}
$$

Notice that $\mathbb{P}(\gamma, -) = 0$ breaks the internal Lorentz gauge. From the above equations, we can parameterize the variables $H^{jk}_a$ and $H^{0i}_a$ purely down to $\mathbb{P}(\gamma, +)^i_{ab}$ and variables $N, N^i$, which we can see explicitly by solving for them:

$$
\begin{aligned}
&\mathbb{H}\left(\frac{1}{\gamma}, -\right)^i_a - \mathbb{H}(\gamma, -)^i_a = \left(\frac{1}{2\gamma} - \frac{\gamma}{2}\right) H^{0i}_a = \frac{1-\gamma^2}{2\gamma} H^{0i}_a \quad &\text{(i)}\\
&\gamma^2 \mathbb{H}(\gamma, -)^i_a - \mathbb{H}\left(\frac{1}{\gamma}, -\right)^i_a = \frac{\gamma^2 - 1}{4} \epsilon^i_{\ jk} H^{jk}_a \quad &\text{(ii)}
\end{aligned}
$$

Plugging in our findings for $\mathbb{H}(\gamma, -)$ and $\mathbb{H}\left(1/\gamma, -\right)$, we get:

$$
\begin{aligned}
\text{(i)} &\implies \frac{1 - \gamma^2}{2} N^b \mathbb{P}(\gamma, +)^i_{ba} + N \frac{1 - \gamma^2}{2\gamma^2} e^i_a = \frac{1-\gamma^2}{2\gamma} H^{0i}_a \\
&\implies H^{0i}_a = \frac{N}{\gamma} e^i_a + \gamma N^b \mathbb{P}(\gamma, +)^i_{ba} \\
\text{(ii)} &\implies - N \frac{1 - \gamma^2}{2} e^i_a - \frac{1 - \gamma^2}{2} N^b \mathbb{P}(\gamma, +)^i_{ba} = \frac{\gamma^2 - 1}{4} \epsilon^i_{\ jk} H^{jk}_a \\
&\implies \frac{1}{2} \epsilon^{i}_{\ jk} H^{jk}_a = N e^i_a + N^b \mathbb{P}(\gamma, +)^i_{ba}
\end{aligned}
$$

Thus, we have parameterized the Plebanski constraints in terms of $\mathbb{P}(\gamma, +)^i_{ab}$, a scalar lapse $N$, shift vector $N^a \in T(\Sigma_t)$, and the 3 parameters in the choice of an internal direction $n^I = (1, 0, 0, 0)$. These sum up to $9 + 1 + 3 + 3 = 16$ parameters, which is exactly the 16 degrees of freedom in the co-tetrad $e^I_\mu$.

The constraints are "simple" constraints and they must be conserved in time. By arXiv:0902.3416 [gr-qc], this follows from

$$
\mathbb{C}^{ij} = e^{(j} \dot{\mathbb{P}}(\gamma, -)^{i)} \propto N e^{(j}\wedge \left[\mathrm{d}e^{i)} + \frac{1}{2}\epsilon^{i)}_{\ \ \ lm}\left({}^+\mathbb{A} + {}^-\mathbb{A}\right)^l\wedge e^{m} \right]\approx 0.
$$

Together with $\mathbb{B}^i \approx 0$, these 9 constraints are equivalent to

$$
\mathbb{IV}^i_a = \frac{1}{2} ({}^{+}\mathbb{A} + {}^{-}\mathbb{A})_a^i - \Gamma_a^i(e) = 0.
$$

Where $\Gamma_a^i$ is the spin connection compatible with the triad $e$, i.e., the solution to

$$
\partial_{[a} e_{b]}^i + \epsilon^i_{\ jk} \Gamma_{[a}^j e_{b]}^k = 0.
$$

We can also calculate the Poisson bracket:

$$
\{ \mathbb{IV}^i_a(x), \mathbb{P}(\gamma, -)^j_{bc}(y) \} = \frac{1}{2} \epsilon_{abc} \delta^{ij} \delta(x-y).
$$

Which makes them secondary constraints. By setting them strongly to zero, we define $K_a^i = ({}^{+}\mathbb{A}^i_a - {}^{-}\mathbb{A}^i_a)/2$, such that:

$$
{}^{+}\mathbb{A}_a^i = \Gamma_a^i + \gamma K_a^i =: A_a^i.
$$

The action is then reduced to:

$$
\begin{aligned}
S[\mathbb{P}(\gamma, +), {}^{+}\mathbb{A}, N, N_i]
&= \frac{1}{\kappa} \int dt \int_{\Sigma} \Big( \mathbb{P}(\gamma, +)_i \wedge {}^+\dot{\mathbb{A}}^i \\
&+ N_i \mathcal{D}_{A} \wedge \mathbb{P}(\gamma, +)^i + \epsilon^i_{\ jk} N^j \mathbb{P}(\gamma, +)^k \wedge (\gamma F_{0i} + F^{jk}\epsilon_{ijk}) \\
&+ Ne^i \wedge (\frac{F_{0i}}{\gamma} +  F^{jk}\epsilon_{ijk} )\Big).
\end{aligned}
$$

This is the standard Hamiltonian formulation of general relativity in terms of SU(2) connection variables $A_a^i$. If we now introduce the densitized triad

$$
E^a_i = \gamma \epsilon^{abc} \mathbb{P}(\gamma, +)_{bc}^i,
$$

The Poisson bracket between $E^a_i$ and $A^j_b$ is now

$$
\{E^a_i(x), A^j_b(y)\} = \kappa \delta^a_b \delta^j_i \delta(x-y), \quad \{E^a_i(x), E^b_j(y)\} = \{A^i_a(x), A^j_b(y)\} = 0.
$$

And the action takes the familiar Ashtekar-Barbero form:

$$
\begin{aligned}
S[E, A, N, N^i, N_i] &= \frac{1}{\gamma\kappa} \int dt \int_{\Sigma_t} d^3x \Big[ E^a_i \dot{A}^i_a - N^b V_b(E^a_i, A^i_a) \\
&- N S(E^a_i, A^i_a) - N^i G_i(E^a_i, A^i_a) \Big]
\end{aligned}
$$

With the constraints

$$
\begin{cases}
V_b = E^a_j F^j_{ab} - (1+\gamma^2)K^i_b G_i \\
S = \frac{E^a_i E^b_j}{\sqrt{|\det E|}} \left( \epsilon^{ij}_{\ \ k} F^k_{ab} - 2(1+\gamma^2)K^i_{[a} K^j_{b]} \right) \\
G_i = \mathcal{D}_a E^a_i
\end{cases}
$$

These three constraints completely dictate both the symmetries and the dynamics of General Relativity. Because it is a fully constrained system, its Hamiltonian is simply a linear combination of these constraints. Physically, they can be interpreted as follows:

\subsubsection{1. The Gauss Constraint ($G_i \approx 0$)}

This is the non-Abelian equivalent of Gauss's Law from Yang-Mills theory ($\nabla \cdot \vec{E} = 0$). Here, $E^a_i$ acts as the electric field conjugate to the connection $A_a^i$. It generates local $SU(2)$ internal gauge rotations. By setting $G_i \approx 0$, we ensure that the physical state of the universe is independent of how we locally orient our internal 3D reference frame (the target space of the triads). Picking a different internal $SU(2)$ basis at any point in space doesn't change the physical geometry.

\subsubsection{2. The Vector Constraint ($V_b \approx 0$)}

In the action, this constraint is multiplied by the shift vector $N^b$. It is the generator of spatial diffeomorphisms upon the 3D slice $\Sigma$. Setting $V_b \approx 0$ enforces the principle that coordinates are arbitrary—if you stretch, twist, or remap the spatial coordinate grid on $\Sigma$, the underlying physical geometry remains exactly the same. The $E^a_j F^j_{ab}$ term essentially comes from the Lie derivative of the connection along the spatial slice, while the $- (1+\gamma^2)K^i_b G_i$ part is a proportional correction term.

\subsubsection{3. The Scalar Constraint ($S \approx 0$)}

In the action, this is multiplied by the scalar lapse function $N$. This constraint generates "time evolution," which in General Relativity corresponds to deforming the spatial slice $\Sigma_t$ forward in the normal direction into the 4D spacetime bulk. Because this must also vanish ($S \approx 0$), it embodies the fact that there is no absolute background time in GR. This is the core dynamical equation of General Relativity.



\end{document}
